\section{The whole project on one page}

\begin{kid}
Imagine a calculator that has a \emph{plus} button but no \emph{times} button. You can still
work out $13\times 11$: add 13 to itself 11 times, or use a smarter trick with doubling.
Either way it takes many button presses. Now imagine you have to do this
\textbf{two hundred thousand times} just to recognise one picture of a shoe. That is the
situation of the small chip in this project.

The paper we follow says: ``Let's give the calculator one extra, very cheap button. It doesn't
multiply. It only tells you how big a number roughly is (where its leading 1 is). With that
button, we can get a \emph{nearly} right answer in very few presses.''

This project checks whether the ``nearly right'' answers are good enough for a neural network that
recognises pictures, and how much faster and cheaper it gets.
\end{kid}

\subsection*{The questions we asked}
\begin{enumerate}
  \item \textbf{Accuracy:} if every multiplication in a neural network is a bit wrong, does the
        network still recognise the pictures? (Answered on a GPU in Phase~1.)
  \item \textbf{Speed:} on a real multiplierless CPU with the new LOD instruction, how many clock
        cycles does one picture take with each multiplier? (Answered on a simulated Ibex
        RISC-V core in Phase~2.)
  \item \textbf{Energy:} the LOD instruction makes the chip use slightly more power. Do the saved
        cycles still save energy overall?
  \item \textbf{Can we do even better?} We also tried writing the code more cleverly (``fusing''),
        an exact table-of-squares method that needs no multiplier at all, and skipping zeros.
\end{enumerate}

\subsection*{The answers in one table}

\begin{table}[H]
\centering
\small
\begin{tabular}{lcccc}
\toprule
 & \textbf{Accuracy change} & \multicolumn{3}{c}{\textbf{Speed on Ibex vs.\ exact (fused code)}}\\
\cmidrule(lr){3-5}
 & (GPU, int8) & MLP & LeNet-5 & Energy vs.\ exact\\
\midrule
Exact shift-and-add & reference & 1.00$\times$ & 1.00$\times$ & 1.00\\
RoBA-Dn1 (needs LOD) & at most $-0.75$ points & 1.91--1.95$\times$ & 2.25--2.36$\times$ & 0.44--0.55\\
Square table (no LOD) & none (exact) & 3.49--3.59$\times$ & 3.03--3.17$\times$ & 0.28--0.33\\
Table + raw pixels + zero skip & none (exact) & 5.0--10.5$\times$ & 2.9--3.3$\times$ & ---\\
\bottomrule
\end{tabular}
\caption{The main results. Ranges cover the two datasets (MNIST and Fashion-MNIST). ``Points''
means percentage points of test accuracy, so 90.95\% $\to$ 90.20\% is $-0.75$ points.}
\end{table}

\subsection*{How this document is organised}
\begin{itemize}
  \item \textbf{Section~\ref{sec:basics}} starts from zero: bits, shifting, how a computer
        multiplies with only additions, what a CPU and a clock cycle are, and what the LOD
        instruction does.
  \item \textbf{Section~\ref{sec:mult}} explains all eight multipliers with the same small example,
        $13\times 11$, so you can follow every step by hand.
  \item \textbf{Section~\ref{sec:nn}} explains neural networks, our two networks, and how decimals
        are turned into whole numbers (quantization).
  \item \textbf{Section~\ref{sec:gpu}} shows the GPU accuracy results.
  \item \textbf{Section~\ref{sec:ibex}} explains the Ibex CPU, how we added LOD to it, and how a
        trained network ends up running on the simulated chip.
  \item \textbf{Section~\ref{sec:ibexres}} shows every Ibex result: speed, where the time goes,
        the faster code, the table of squares, zero skipping, and energy.
  \item \textbf{Section~\ref{sec:wrap}} sums up, lists the limits, and has a glossary.
  \item \textbf{The appendix} walks through every code file and gives the commands that rerun
        everything.
\end{itemize}

Throughout the document you will see coloured boxes:
\begin{itemize}
  \item \textbf{Explained simply} (orange): the idea in everyday words.
  \item \textbf{Worked example} (blue): small numbers you can check with pencil and paper.
  \item \textbf{Key idea} (green): the one thing to remember from that part.
  \item \textbf{Be careful} (red): a common misunderstanding, or a limit of what we measured.
  \item \textbf{Where this is in the code} (grey): which file and function does what was just
        explained.
\end{itemize}
\newpage
